\section{Performance Beyond Proportional Compute}
\label{sec:evidence}

The promise of allocation intelligence must be visible from outside the
system. Fusion should not require a customer to choose between frontier-level
capability and practical deployment efficiency. It should move the frontier
itself.

\subsection{A performance--compute frontier across three capability regimes}

Our current evaluation spans scientific reasoning, policy-constrained tool use,
and agentic coding. These workloads expose different failure modes, but ask the
same systems question: how much capability can be delivered for the compute
consumed? Table~\ref{tab:pareto} reports both score and the cost in US dollars
of completing each full benchmark run.

\begin{table}[H]
\centering
\caption{Reported score and full benchmark-run cost. Lower cost and higher
score are better.}
\label{tab:pareto}
\scriptsize
\begin{tabular}{@{}l
  r@{\hspace{1.2em}}r
  @{\hspace{2.0em}}r@{\hspace{1.2em}}r
  @{\hspace{2.0em}}r@{\hspace{1.2em}}r@{}}
\toprule
& \multicolumn{2}{c}{GPQA Diamond}
& \multicolumn{2}{c}{$\tau^3$-Banking}
& \multicolumn{2}{c}{Terminal-Bench v2.1} \\
\cmidrule(lr){2-3}\cmidrule(lr){4-5}\cmidrule(lr){6-7}
System & Score & Cost & Score & Cost & Score & Cost \\
\midrule
\textbf{Fusion Max} & \textbf{94.9\%} & \textbf{13}
  & \textbf{49.5\%} & \textbf{145} & 86.5\% & \textbf{74} \\
GPT 5.6 Sol (max) & 94.1\% & 22 & 44.3\% & 251
  & \textbf{88.0\%} & 86 \\
Claude Fable 5 (fallback) & 92.6\% & 45 & 38.1\% & 322
  & 84.6\% & 299 \\
\addlinespace
\textbf{Fusion Code} & \textbf{93.5\%} & \textbf{7}
  & \textbf{48.5\%} & 159 & \textbf{89.9\%} & 71 \\
Claude Opus 5 (max) & 93.2\% & 18 & 42.1\% & 193
  & 89.1\% & 215 \\
Kimi K3 & \textbf{93.5\%} & 30 & 46.0\% & \textbf{122}
  & 85.0\% & \textbf{56} \\
\bottomrule
\end{tabular}
\end{table}

Fusion Max reduces the aggregate cost of this fixed three-benchmark portfolio
from \$359 to \$232 relative to GPT 5.6 Sol, a 35.4\% reduction. It scores
higher on GPQA Diamond and $\tau^3$-Banking and remains within 1.5 percentage
points on Terminal-Bench. Against Claude Fable 5, Fusion Max scores higher on
all three workloads while reducing each full-run cost by 55--75\%; aggregate
cost falls from \$666 to \$232.

Fusion Code presents a second operating point. It meets or exceeds the reported
score of Claude Opus 5 on all three workloads while reducing full-run cost by
18--67\%; aggregate cost falls from \$426 to \$237. Kimi K3 illustrates a
different trade-off: Fusion Code matches or improves its score on all three
workloads, but spends more on the Banking and Terminal-Bench runs. We therefore
present that comparison as a capability gain, not as a universal cost
reduction.

\evidencepending{before publication, attach immutable run identifiers and
confirm model snapshots, pricing date, task scope, scorer and protocol
versions, attempts and retries, fallback behavior, cache treatment, and the
boundary of included internal usage.}

The result is a Pareto claim, not a discount claim. Fusion retains stronger
judgment where it is valuable and removes the assumption that every
intermediate unit of reasoning must be purchased at the same capability tier.
A system that is only cheaper because it silently accepts worse outcomes has
not demonstrated allocation intelligence.

\subsection{Why asymmetry changes the economics of capability}

The cost profile follows directly from Fusion's architecture. A frontier-heavy
orchestration design assembles its worker pool from several of the most capable
general-purpose models. This can increase quality, but it also duplicates the
most expensive class of computation across planning, generation, critique, and
verification. Fusion instead composes a wider capability spectrum. It asks
where frontier intelligence is decisive and where a specialized or efficient
source can add more information per unit of compute.

Fugu provides a useful external operating point. It uses learned orchestration
across a frontier-heavy worker pool, while Fusion uses task-state understanding
to allocate across an asymmetric supply of models. On the benchmark comparison
available to us, the two systems reach broadly comparable performance, while
Fugu's reported benchmark cost is roughly an order of magnitude higher. The
difference is not incidental pricing. It reflects two system designs: one
primarily orchestrates frontier intelligence; the other learns where frontier
intelligence is necessary.

\evidencepending{insert the audited Fugu case study, with matched task subset,
model and price dates, scaffold differences, attempts, usage boundary, and a
clear distinction between controlled and directional comparisons.}

The point is not that an asymmetric pool is inherently superior. A weak source
that adds misleading or redundant evidence is not efficient at any price
\citep{li2025rethinkingmoa}. The advantage appears only when the system has the
capability understanding and task-state awareness to use asymmetry
selectively. That allocation layer---not the mere presence of smaller
models---is the source of the frontier shift.

\subsection{Effectiveness that survives the workflow}

One-shot benchmarks capture only part of Fusion's value. In an agentic system,
the demand for intelligence changes after every observation and action. A
design that improves a single answer can still fail over a trajectory by
repeating work, losing state, acting on incompatible plans, or declaring
success before verification.

Fusion is designed so that its benefit persists across the workflow. The same
canonical state, uncertainty model, and authority boundary apply whether the
system is investigating, editing, testing, recovering from failure, or
committing a result. The relevant outcome becomes successful task completion,
not the quality of an isolated response. The relevant cost becomes compute per
successful trajectory, not price per API call.

Our long-horizon evaluation therefore measures end-to-end completion, total
cost per success, frontier calls per trajectory, state consistency, recovery
after failed actions, premature completion, and degradation as workflow length
increases. This is where allocation intelligence becomes a system capability
rather than a routing optimization.

\subsection{A transparent standard of evidence}

Public comparisons should expose the full operating point: task set, model and
policy versions, scorer, attempts, timeout, concurrency, cache treatment,
latency, and all internal model usage. Client-visible tokens are not the total
cost of orchestration. Infrastructure and scorer failures must remain separate
from system failures. Fumo Lab will report these boundaries explicitly,
especially where a fast-moving external system cannot be reproduced under a
perfectly matched protocol.
